% appP_register.tex -- generated by writing/tools/blueprint_gen.py 1.0 (level 3) on 2026-10-14 from sources 20021204b9; do not edit
\providecommand{\bpNum}[1]{\hyperref[#1]{\ref*{#1}}}%
\providecommand{\bpNumA}[1]{\hyperref[#1]{\textcolor{orange!65!black}{\ref*{#1}}}}%
\providecommand{\bpPage}[1]{\pageref*{#1}}%
\providecommand{\bpLab}[1]{{\ttfamily #1}}%
\providecommand{\bpBrk}{\allowbreak}%
\providecommand{\bpT}[2]{\facttagbox[#1]{\scriptsize #2}}%
\providecommand{\bpU}[1]{{\color{gray!70!black}#1}}%

\section{假设登记表}\label{sec:appP-register}

本表由各单元的 \texttt{assumption} 环境与假设登记表自动生成，按所属单元的书序排列；“被使用于”一列由各陈述的 \texttt{\string\uses} 数据计算（直接使用者的编号），与环境中「被使用于」字段的比较见检查 C-P10。
\begingroup\scriptsize\setlength{\tabcolsep}{2.5pt}\renewcommand{\arraystretch}{1.06}
\begin{longtable}{@{}>{\raggedright\arraybackslash}p{\dimexpr(\textwidth-12\tabcolsep-0mm)*151/1000\relax}>{\raggedright\arraybackslash}p{\dimexpr(\textwidth-12\tabcolsep-0mm)*75/1000\relax}>{\raggedright\arraybackslash}p{\dimexpr(\textwidth-12\tabcolsep-0mm)*158/1000\relax}>{\raggedright\arraybackslash}p{\dimexpr(\textwidth-12\tabcolsep-0mm)*49/1000\relax}>{\raggedright\arraybackslash}p{\dimexpr(\textwidth-12\tabcolsep-0mm)*138/1000\relax}>{\raggedright\arraybackslash}p{\dimexpr(\textwidth-12\tabcolsep-0mm)*303/1000\relax}>{\raggedright\arraybackslash}p{\dimexpr(\textwidth-12\tabcolsep-0mm)*122/1000\relax}@{}}
\caption{编号假设登记表（由 \texttt{assumption} 环境生成）}\label{tab:appP-register}\\
\toprule
\textbf{编号／标签} & \textbf{单元} & \textbf{内容} & \textbf{标签} & \textbf{地位} & \textbf{被使用于} & \textbf{解除途径} \\
\midrule
\endfirsthead
\caption[]{编号假设登记表（由 \texttt{assumption} 环境生成）（续）}\\
\toprule
\textbf{编号／标签} & \textbf{单元} & \textbf{内容} & \textbf{标签} & \textbf{地位} & \textbf{被使用于} & \textbf{解除途径} \\
\midrule
\endhead
\midrule
\multicolumn{7}{r@{}}{\footnotesize（续下页）}\\
\endfoot
\bottomrule
\endlastfoot
\bpNumA{asm:intro-wightman}\newline \bpLab{\detokenize{asm:}\bpBrk \detokenize{intro-}\bpBrk \detokenize{wightman}} & \mbox{第~1~章} & QCD 的 Wightman 结构 (W1)--(W4) & \bpT{Plum}{C} & 全局非微扰假设。测试状态句。 & \hyperref[asm:intro-composite]{\textcolor{orange!65!black}{1.10}}, \bpBrk \hyperref[asm:intro-np-ope]{\textcolor{orange!65!black}{11}}, \bpBrk \hyperref[asm:intro-growth]{\textcolor{orange!65!black}{13}}, \bpBrk \hyperref[def:intro-MA]{15}；\bpBrk \hyperref[prop:prelim-optical]{3.24}；\bpBrk \hyperref[prop:ee-R-ImPi]{6.2}, \bpBrk \hyperref[prop:ee-analytic]{4}, \bpBrk \hyperref[prop:ee-smearing]{23}；\bpBrk \hyperref[prop:dis-W-structure]{7.2}, \bpBrk \hyperref[prop:dis-ope-rg]{17}, \bpBrk \hyperref[prop:dis-dispersion-np]{19}, \bpBrk \hyperref[thm:dis-fixed-n]{22}--\hyperref[thm:dis-moments]{25}, \bpBrk \hyperref[prop:dis-twist4]{33}, \bpBrk \hyperref[prop:dis-tmc]{34}, \bpBrk \hyperref[thm:dis-bjorken]{36}；\bpBrk \hyperref[prop:pdf-invariance]{8.4}, \bpBrk \hyperref[prop:pdf-support-hadron]{7}, \bpBrk \hyperref[thm:pdf-moments-local]{13}, \bpBrk \hyperref[thm:pdf-dis-xspace]{22}, \bpBrk \hyperref[prop:pdf-sum-rules]{23}, \bpBrk \hyperref[prop:pdf-timelike-dglap]{29}；\bpBrk \hyperref[prop:coll-photon-sf]{11.10}, \bpBrk \hyperref[prop:coll-luxqed]{13}, \bpBrk \hyperref[thm:coll-g1-moments]{15}, \bpBrk \hyperref[prop:coll-transversity-evolution]{21}；\bpBrk \hyperref[thm:tmd-sivers-sign]{14.20}；\bpBrk \hyperref[prop:excl-erbl]{19.2}, \bpBrk \hyperref[prop:excl-polynomiality]{10}；\bpBrk \hyperref[thm:heavy-inclusive-b]{20.5}；\bpBrk \hyperref[thm:lamet-exp-decay]{22.32}；\bpBrk \hyperref[lem:appM-forward-tube]{M.54}, \bpBrk \hyperref[lem:appM-positivity]{57} & 测试解除途径。 \\
\bpNumA{asm:intro-composite}\newline \bpLab{\detokenize{asm:}\bpBrk \detokenize{intro-}\bpBrk \detokenize{composite}} & \mbox{第~1~章} & 规范不变复合算符与光线算符的存在与重整化群 & \bpT{Plum}{C} & 全局非微扰假设。测试状态句。 & \hyperref[asm:intro-np-ope]{\textcolor{orange!65!black}{1.11}}, \bpBrk \hyperref[def:intro-MA]{15}；\bpBrk \hyperref[prop:dis-ope-rg]{7.17}, \bpBrk \hyperref[thm:dis-moments-np]{23}--\hyperref[thm:dis-moments]{25}, \bpBrk \hyperref[prop:dis-twist4]{33}, \bpBrk \hyperref[prop:dis-tmc]{34}, \bpBrk \hyperref[thm:dis-bjorken]{36}；\bpBrk \hyperref[prop:pdf-invariance]{8.4}, \bpBrk \hyperref[prop:pdf-support-hadron]{7}, \bpBrk \hyperref[thm:pdf-moments-local]{13}, \bpBrk \hyperref[thm:pdf-dglap]{14}, \bpBrk \hyperref[thm:pdf-dis-xspace]{22}, \bpBrk \hyperref[prop:pdf-sum-rules]{23}, \bpBrk \hyperref[prop:pdf-timelike-dglap]{29}；\bpBrk \hyperref[prop:coll-photon-sf]{11.10}, \bpBrk \hyperref[prop:coll-luxqed]{13}, \bpBrk \hyperref[thm:coll-g1-moments]{15}, \bpBrk \hyperref[prop:coll-transversity-evolution]{21}；\bpBrk \hyperref[prop:tmd-small-b]{14.18}, \bpBrk \hyperref[thm:tmd-sivers-sign]{20}；\bpBrk \hyperref[prop:excl-erbl]{19.2}, \bpBrk \hyperref[prop:excl-polynomiality]{10}；\bpBrk \hyperref[thm:heavy-inclusive-b]{20.5}；\bpBrk \hyperref[lem:appM-norm-bound]{M.59} & 测试解除途径。 \\
\bpNumA{asm:intro-np-ope}\newline \bpLab{\detokenize{asm:}\bpBrk \detokenize{intro-}\bpBrk \detokenize{np-}\bpBrk \detokenize{ope}} & \mbox{第~1~章} & 非微扰算符乘积展开 (O1)(O2) & \bpT{Plum}{C} & 全局非微扰假设。测试状态句。 & \hyperref[asm:intro-af-coeff]{\textcolor{orange!65!black}{1.12}}, \bpBrk \hyperref[def:intro-MA]{15}；\bpBrk \hyperref[lem:dis-ope-segment]{7.21}--\hyperref[thm:dis-moments]{25}, \bpBrk \hyperref[prop:dis-twist4]{33}, \bpBrk \hyperref[prop:dis-tmc]{34}, \bpBrk \hyperref[thm:dis-bjorken]{36}；\bpBrk \hyperref[thm:pdf-dis-xspace]{8.22}；\bpBrk \hyperref[prop:coll-photon-sf]{11.10}, \bpBrk \hyperref[thm:coll-g1-moments]{15} & 测试解除途径。 \\
\bpNumA{asm:intro-af-coeff}\newline \bpLab{\detokenize{asm:}\bpBrk \detokenize{intro-}\bpBrk \detokenize{af-}\bpBrk \detokenize{coeff}} & \mbox{第~1~章} & 系数的渐近自由行为 & \bpT{Plum}{C} & 全局非微扰假设。测试状态句。 & \hyperref[def:intro-MA]{1.15}；\bpBrk \hyperref[thm:dis-moments-np]{7.23}--\hyperref[thm:dis-moments]{25}, \bpBrk \hyperref[prop:dis-twist4]{33}, \bpBrk \hyperref[prop:dis-tmc]{34}, \bpBrk \hyperref[thm:dis-bjorken]{36}；\bpBrk \hyperref[thm:pdf-dis-xspace]{8.22}；\bpBrk \hyperref[prop:coll-photon-sf]{11.10}, \bpBrk \hyperref[thm:coll-g1-moments]{15}；\bpBrk \hyperref[thm:heavy-inclusive-b]{20.5}；\bpBrk \hyperref[lem:appM-norm-bound]{M.59} & 测试解除途径。 \\
\bpNumA{asm:intro-growth}\newline \bpLab{\detokenize{asm:}\bpBrk \detokenize{intro-}\bpBrk \detokenize{growth}} & \mbox{第~1~章} & 前向振幅的增长条件 & \bpT{Plum}{C} & 全局非微扰假设。测试状态句。 & \hyperref[def:intro-MA]{1.15}；\bpBrk \hyperref[thm:dis-moments-np]{7.23}, \bpBrk \hyperref[thm:dis-moments]{25}, \bpBrk \hyperref[thm:dis-bjorken]{36}；\bpBrk \hyperref[thm:pdf-dis-xspace]{8.22} & 测试解除途径。 \\
\bpNumA{asm:intro-bridge}\newline \bpLab{\detokenize{asm:}\bpBrk \detokenize{intro-}\bpBrk \detokenize{bridge}} & \mbox{第~1~章} & 从部分子微扰论到强子的桥接假设 & \bpT{Plum}{C} & 全局非微扰假设。测试状态句。 & \hyperref[def:intro-MA]{1.15}；\bpBrk \hyperref[thm:method-template]{9.42}；\bpBrk \hyperref[thm:disall-factorization]{10.11}, \bpBrk \hyperref[thm:disall-heavy]{18}, \bpBrk \hyperref[prop:disall-nuclear]{20}；\bpBrk \hyperref[thm:coll-sia]{11.5}, \bpBrk \hyperref[thm:coll-sidis]{8}, \bpBrk \hyperref[thm:coll-direct-photo]{9}；\bpBrk \hyperref[thm:dy-factorization]{12.17}, \bpBrk \hyperref[thm:dy-colour-singlet]{25}, \bpBrk \hyperref[thm:dy-polarized]{26}, \bpBrk \hyperref[thm:dy-higher-twist]{29}；\bpBrk \hyperref[thm:tmd-factorization]{14.15}, \bpBrk \hyperref[prop:tmd-small-b]{18}, \bpBrk \hyperref[cor:tmd-sivers-observable]{21}, \bpBrk \hyperref[prop:tmd-collins-universality]{22}；\bpBrk \hyperref[thm:scet-thrust]{15.16}；\bpBrk \hyperref[thm:jets-ngl-factorization]{16.10}；\bpBrk \hyperref[thm:break-ddis]{18.14}, \bpBrk \hyperref[thm:break-double-dy]{17}；\bpBrk \hyperref[thm:excl-pion-ff]{19.4}, \bpBrk \hyperref[thm:excl-transition]{7}, \bpBrk \hyperref[thm:excl-dvcs]{11}, \bpBrk \hyperref[thm:excl-meson-L]{13}；\bpBrk \hyperref[thm:heavy-bdpi]{20.1}--\hyperref[prop:heavy-qcdf-lp]{3}, \bpBrk \hyperref[thm:heavy-inclusive-b]{5}, \bpBrk \hyperref[thm:heavy-nrqcd-decay]{7}, \bpBrk \hyperref[prop:heavy-quarkonium-fragmentation]{11}；\bpBrk \hyperref[prop:power-dis-twist4-x]{21.1} & 测试解除途径。 \\
\bpNumA{asm:lamet-aux-renorm}\newline \bpLab{\detokenize{asm:}\bpBrk \detokenize{lamet-}\bpBrk \detokenize{aux-}\bpBrk \detokenize{renorm}} & \mbox{第~22~章} & 含类空辅助静态色源的 QCD 的可重整性 (AX1)--(AX3) & \bpT{RoyalBlue}{A} & 测试状态句。 & \hyperref[prop:lattice-renormalizability]{22.9}, \bpBrk \hyperref[asm:lamet-ope]{\textcolor{orange!65!black}{19}}, \bpBrk \hyperref[asm:lamet-gap]{\textcolor{orange!65!black}{31}}, \bpBrk \hyperref[conj:lamet-decay-qcd]{33}, \bpBrk \hyperref[asm:lamet-staple]{\textcolor{orange!65!black}{53}}；\bpBrk \hyperref[thm:yzliu-soft-hqet]{23.4} & 测试解除途径。 \\
\bpNumA{asm:lamet-ope}\newline \bpLab{\detokenize{asm:}\bpBrk \detokenize{lamet-}\bpBrk \detokenize{ope}} & \mbox{第~22~章} & 类空 Wilson 线算符的短距离 OPE（重整化微扰论） (OP1)(OP2) & \bpT{RoyalBlue}{A} & 测试状态句。 & \hyperref[thm:lamet-coordinate]{22.20}, \bpBrk \hyperref[thm:lamet-parton-momentum]{21}, \bpBrk \hyperref[prop:lamet-tmc]{38}, \bpBrk \hyperref[prop:lamet-sdf]{48}, \bpBrk \hyperref[thm:lamet-quasi-gpd]{51} & 测试解除途径。 \\
\bpNumA{asm:lamet-gap}\newline \bpLab{\detokenize{asm:}\bpBrk \detokenize{lamet-}\bpBrk \detokenize{gap}} & \mbox{第~22~章} & 谱隙与边界值正则性 (GP1)--(GP3) & \bpT{Plum}{C} & 测试状态句。 & \hyperref[thm:lamet-exp-decay]{22.32}, \bpBrk \hyperref[conj:lamet-decay-qcd]{33} & 测试解除途径。 \\
\bpNumA{asm:lamet-uniform}\newline \bpLab{\detokenize{asm:}\bpBrk \detokenize{lamet-}\bpBrk \detokenize{uniform}} & \mbox{第~22~章} & 扭度展开对 Ioffe 时间的一致性 & \bpT{Plum}{C} & 测试状态句。 & \hyperref[thm:lamet-hadronic]{22.46} & 测试解除途径。 \\
\bpNumA{asm:lamet-twist4}\newline \bpLab{\detokenize{asm:}\bpBrk \detokenize{lamet-}\bpBrk \detokenize{twist4}} & \mbox{第~22~章} & 领头幂次修正的结构（扭度 4） (TW1)(TW2) & \bpT{Plum}{C} & 测试状态句。 & \hyperref[asm:lamet-large-z]{\textcolor{orange!65!black}{22.45}}, \bpBrk \hyperref[thm:lamet-hadronic]{46} & 测试解除途径。 \\
\bpNumA{asm:lamet-large-z}\newline \bpLab{\detokenize{asm:}\bpBrk \detokenize{lamet-}\bpBrk \detokenize{large-}\bpBrk \detokenize{z}} & \mbox{第~22~章} & 大距离区域的控制 (LZ1)(LZ2) & \bpT{Plum}{C} & 测试状态句。 & \hyperref[thm:lamet-hadronic]{22.46} & 测试解除途径。 \\
\bpNumA{asm:lamet-staple}\newline \bpLab{\detokenize{asm:}\bpBrk \detokenize{lamet-}\bpBrk \detokenize{staple}} & \mbox{第~22~章} & 订书钉形类空 Wilson 线算符的重整化与 $L\to\infty$ 极限 (ST1)--(ST3) & \bpT{RoyalBlue}{A} & 测试状态句。 & \hyperref[def:lamet-quasi-tmd]{22.54}, \bpBrk \hyperref[thm:lamet-quasi-tmd]{59}；\bpBrk \hyperref[prop:yzliu-qtmd-powercount]{23.23} & 测试解除途径。 \\
\bpNumA{asm:lamet-soft-rapidity}\newline \bpLab{\detokenize{asm:}\bpBrk \detokenize{lamet-}\bpBrk \detokenize{soft-}\bpBrk \detokenize{rapidity}} & \mbox{第~22~章} & 离光锥软函数的大快度结构 (SR) & \bpT{RoyalBlue}{A} & 测试状态句。 & \hyperref[def:lamet-reduced-soft]{22.57}, \bpBrk \hyperref[thm:lamet-quasi-tmd]{59} & 测试解除途径。 \\
\bpNumA{asm:lamet-order-limits}\newline \bpLab{\detokenize{asm:}\bpBrk \detokenize{lamet-}\bpBrk \detokenize{order-}\bpBrk \detokenize{limits}} & \mbox{第~22~章} & 极限次序 (OL1)(OL2) & \bpT{RoyalBlue}{A} & 测试状态句。 & \hyperref[thm:lamet-quasi-tmd]{22.59} & 测试解除途径。 \\
\bpNumA{asm:yzliu-heavy-pair}\newline \bpLab{\detokenize{asm:}\bpBrk \detokenize{yzliu-}\bpBrk \detokenize{heavy-}\bpBrk \detokenize{pair}} & \mbox{第~23~章} & 静态色单态 $Q\bar Q$ 对的谱假设 (HP1)--(HP3) & \bpT{Plum}{C} & 测试状态句。 & \hyperref[cor:yzliu-soft-form-factor]{23.6} & 测试解除途径。 \\
\bpNumA{asm:yzliu-wightman-pt}\newline \bpLab{\detokenize{asm:}\bpBrk \detokenize{yzliu-}\bpBrk \detokenize{wightman-}\bpBrk \detokenize{pt}} & \mbox{第~23~章} & 维数正规化微扰论中的 Wightman 结构 (WP1)(WP2) & \bpT{RoyalBlue}{A} & 测试状态句。 & \hyperref[thm:yzliu-analyticity]{23.9}, \bpBrk \hyperref[thm:yzliu-time-ordering]{20}, \bpBrk \hyperref[conj:yzliu-complex-wilson]{22}, \bpBrk \hyperref[thm:yzliu-jet-continuation]{31} & 测试解除途径。 \\
\bpNumA{asm:yzliu-lightcone-uniform}\newline \bpLab{\detokenize{asm:}\bpBrk \detokenize{yzliu-}\bpBrk \detokenize{lightcone-}\bpBrk \detokenize{uniform}} & \mbox{第~23~章} & 光锥极限对 ${\rm Im}\,\theta$ 的一致性 (UL) & \bpT{BurntOrange}{P} & 测试状态句。 & \hyperref[prop:yzliu-universality-classes]{23.16} & 测试解除途径。 \\
\bpNumA{asm:front-mor-complete}\newline \bpLab{\detokenize{asm:}\bpBrk \detokenize{front-}\bpBrk \detokenize{mor-}\bpBrk \detokenize{complete}} & \mbox{第~24~章} & 区域集合的完备性 (MR) & \bpT{Plum}{C} & 测试状态句。 & \hyperref[prop:front-commuting]{24.25} & 测试解除途径。 \\
\bpNumA{asm:front-scet-modes}\newline \bpLab{\detokenize{asm:}\bpBrk \detokenize{front-}\bpBrk \detokenize{scet-}\bpBrk \detokenize{modes}} & \mbox{第~24~章} & SCET 的模式完备性与 Glauber 模式的缺席 (SM1)(SM2) & \bpT{RoyalBlue}{A} & 测试状态句。 & \hyperref[thm:front-hgg]{24.40}, \bpBrk \hyperref[prop:front-ggh]{42}, \bpBrk \hyperref[prop:front-nlp-ll]{44}, \bpBrk \hyperref[prop:front-nlp-dis]{45}, \bpBrk \hyperref[prop:front-sidis-nlp]{49} & 测试解除途径。 \\
\bpNumA{asm:appM-remainder}\newline \bpLab{\detokenize{asm:}\bpBrk \detokenize{appM-}\bpBrk \detokenize{remainder}} & \mbox{附录~M} & 冻结 $q$ 双局域正规乘积的动量空间余项估计 & \bpT{RoyalBlue}{A} & 测试状态句。 & \hyperref[thm:appM-scaling]{M.25}, \bpBrk \hyperref[cor:appM-coordinate]{28} & 测试解除途径。 \\
\bpNumA{asm:appM-forward}\newline \bpLab{\detokenize{asm:}\bpBrk \detokenize{appM-}\bpBrk \detokenize{forward}} & \mbox{附录~M} & 前向在壳无质量部分子矩阵元的 OPE 余项 & \bpT{RoyalBlue}{A} & 测试状态句。 & \hyperref[prop:dis-ope-forward]{7.13} & 测试解除途径。 \\
\bpNumA{asm:appM-brst-ope}\newline \bpLab{\detokenize{asm:}\bpBrk \detokenize{appM-}\bpBrk \detokenize{brst-}\bpBrk \detokenize{ope}} & \mbox{附录~M} & OPE 的 BRST 结构：带一个复合旁观者的欧氏展开 (Eε) (B2)(B1) & \bpT{RoyalBlue}{A} & 测试状态句。 & \hyperref[prop:dis-ope-gi]{7.12}；\bpBrk \hyperref[thm:appM-brst-reduction]{M.41} & 测试解除途径。 \\
\bpNumA{asm:appM-oswi}\newline \bpLab{\detokenize{asm:}\bpBrk \detokenize{appM-}\bpBrk \detokenize{oswi}} & \mbox{附录~M} & 前向在壳胶子振幅的 Ward 恒等式 & \bpT{RoyalBlue}{A} & 测试状态句。 & \hyperref[lem:appM-gluon-regularity]{M.51} & 测试解除途径。 \\
\bpNumA{asm:appN-contour}\newline \bpLab{\detokenize{asm:}\bpBrk \detokenize{appN-}\bpBrk \detokenize{contour}} & \mbox{附录~N} & (GA1) 围道 & \bpT{RoyalBlue}{A} & 测试状态句。 & \hyperref[cor:method-removing-eps]{9.32}, \bpBrk \hyperref[thm:method-subtraction]{33}；\bpBrk \hyperref[prop:disall-two-routes]{10.2}, \bpBrk \hyperref[thm:disall-renorm-compat]{8}；\bpBrk \hyperref[thm:appN-subtraction]{N.11}, \bpBrk \hyperref[prop:appN-sia]{34}, \bpBrk \hyperref[thm:appN-structure]{43} & 测试解除途径。 \\
\bpNumA{asm:appN-envelope}\newline \bpLab{\detokenize{asm:}\bpBrk \detokenize{appN-}\bpBrk \detokenize{envelope}} & \mbox{附录~N} & (GA3) 幂次计数包络 & \bpT{RoyalBlue}{A} & 测试状态句。 & \hyperref[lem:method-local-accuracy]{9.12}, \bpBrk \hyperref[thm:method-template]{42}；\bpBrk \hyperref[prop:disall-two-routes]{10.2}, \bpBrk \hyperref[thm:disall-factorization]{11}；\bpBrk \hyperref[rem:appE-actual-bound]{E.2}；\bpBrk \hyperref[lem:appN-error-table]{N.27}, \bpBrk \hyperref[prop:appN-sia]{34}, \bpBrk \hyperref[thm:appN-structure]{43}；\bpBrk \hyperref[prop:appO-one-sided]{O.54}, \bpBrk \hyperref[thm:appO-main]{68}--\hyperref[prop:appO-tmd]{70} & 测试解除途径。 \\
\bpNumA{asm:appN-separation}\newline \bpLab{\detokenize{asm:}\bpBrk \detokenize{appN-}\bpBrk \detokenize{separation}} & \mbox{附录~N} & (GA4) 分离性 & \bpT{RoyalBlue}{A} & 测试状态句。 & \hyperref[lem:method-local-accuracy]{9.12}, \bpBrk \hyperref[thm:method-template]{42}；\bpBrk \hyperref[prop:disall-two-routes]{10.2}, \bpBrk \hyperref[thm:disall-factorization]{11}；\bpBrk \hyperref[lem:appN-error-table]{N.27}, \bpBrk \hyperref[prop:appN-sia]{34}, \bpBrk \hyperref[prop:appN-tmd-window]{41}, \bpBrk \hyperref[thm:appN-structure]{43}；\bpBrk \hyperref[prop:appO-one-sided]{O.54}, \bpBrk \hyperref[thm:appO-main]{68}--\hyperref[prop:appO-tmd]{70} & 测试解除途径。 \\
\bpNumA{asm:appN-lce}\newline \bpLab{\detokenize{asm:}\bpBrk \detokenize{appN-}\bpBrk \detokenize{lce}} & \mbox{附录~N} & 光锥序表示与协变值的等同（LCE） (M1)(M2) & \bpT{RoyalBlue}{A} & 测试状态句。 & \hyperref[cor:appN-cs3-fibered]{N.18} & 测试解除途径。 \\
\bpNumA{asm:appN-values}\newline \bpLab{\detokenize{asm:}\bpBrk \detokenize{appN-}\bpBrk \detokenize{values}} & \mbox{附录~N} & 发散链项的取值规则 (V1)(V2) & 无 & 测试状态句。 & \hyperref[thm:method-template]{9.42}；\bpBrk \hyperref[prop:appN-thrust-zerobin]{N.35}, \bpBrk \hyperref[prop:appN-thrust]{36}, \bpBrk \hyperref[thm:appN-structure]{43}, \bpBrk \hyperref[prop:appN-triangle-exact]{44}, \bpBrk \hyperref[prop:appN-twoloop-remainder]{51} & 测试解除途径。 \\
\bpNumA{asm:appN-collins}\newline \bpLab{\detokenize{asm:}\bpBrk \detokenize{appN-}\bpBrk \detokenize{collins}} & \mbox{附录~N} & 一般图的 Collins 条件 (CS1)(CS3) & \bpT{RoyalBlue}{A} & 测试状态句。 & \hyperref[thm:method-subtraction]{9.33}, \bpBrk \hyperref[thm:method-template]{42}；\bpBrk \hyperref[prop:disall-two-routes]{10.2}, \bpBrk \hyperref[prop:disall-h5]{9}, \bpBrk \hyperref[thm:disall-factorization]{11}；\bpBrk \hyperref[thm:dy-factorization]{12.17}；\bpBrk \hyperref[thm:tmd-factorization]{14.15}；\bpBrk \hyperref[prop:excl-soft-cancel]{19.3}, \bpBrk \hyperref[thm:excl-meson-L]{13}；\bpBrk \hyperref[prop:appN-sia]{N.34}, \bpBrk \hyperref[prop:appN-thrust]{36}, \bpBrk \hyperref[prop:appN-dy]{37}, \bpBrk \hyperref[prop:appN-tmd-window]{41}, \bpBrk \hyperref[thm:appN-structure]{43} & 测试解除途径。 \\
\bpNumA{asm:appN-ward-product}\newline \bpLab{\detokenize{asm:}\bpBrk \detokenize{appN-}\bpBrk \detokenize{ward-}\bpBrk \detokenize{product}} & \mbox{附录~N} & Ward 恒等式的乘积结构与减除的重组 (WI1)(WI2) & \bpT{RoyalBlue}{A} & 测试状态句。 & \hyperref[thm:method-template]{9.42}；\bpBrk \hyperref[prop:disall-factorized-subtraction]{10.5}, \bpBrk \hyperref[thm:disall-factorization]{11}；\bpBrk \hyperref[prop:tmd-zero-bin-soft]{14.11}；\bpBrk \hyperref[prop:excl-soft-cancel]{19.3}；\bpBrk \hyperref[prop:appN-sia]{N.34}, \bpBrk \hyperref[prop:appN-thrust]{36}, \bpBrk \hyperref[prop:appN-dy]{37}, \bpBrk \hyperref[thm:appN-structure]{43}；\bpBrk \hyperref[prop:appO-past-wilson]{O.65}, \bpBrk \hyperref[thm:appO-main]{68}--\hyperref[prop:appO-tmd]{70} & 测试解除途径。 \\
\bpNumA{asm:appN-brst-nonlocal}\newline \bpLab{\detokenize{asm:}\bpBrk \detokenize{appN-}\bpBrk \detokenize{brst-}\bpBrk \detokenize{nonlocal}} & \mbox{附录~N} & 多定域光线算符的 BRST 分类 & \bpT{Plum}{C} & 具名缺口。测试状态句。 & \hyperref[thm:appN-structure]{N.43} & 测试解除途径。 \\
\bpNumA{asm:appN-cutsum}\newline \bpLab{\detokenize{asm:}\bpBrk \detokenize{appN-}\bpBrk \detokenize{cutsum}} & \mbox{附录~N} & 割图求和的公共表示 (CU1)--(CU3) & \bpT{RoyalBlue}{A} & 测试状态句。 & \hyperref[prop:method-fsi]{9.37}, \bpBrk \hyperref[thm:method-template]{42}；\bpBrk \hyperref[prop:disall-two-routes]{10.2}, \bpBrk \hyperref[thm:disall-factorization]{11}；\bpBrk \hyperref[thm:dy-factorization]{12.17}；\bpBrk \hyperref[prop:appN-sia]{N.34}, \bpBrk \hyperref[prop:appN-thrust]{36}, \bpBrk \hyperref[prop:appN-dy]{37}, \bpBrk \hyperref[prop:appN-tmd-window]{41}, \bpBrk \hyperref[thm:appN-structure]{43} & 测试解除途径。 \\
\bpNumA{asm:appN-uv}\newline \bpLab{\detokenize{asm:}\bpBrk \detokenize{appN-}\bpBrk \detokenize{uv}} & \mbox{附录~N} & 被积函数层次的紫外重整化 (U1)(U2) & \bpT{RoyalBlue}{A} & 测试状态句。 & \hyperref[thm:method-subtraction]{9.33}, \bpBrk \hyperref[thm:method-template]{42}；\bpBrk \hyperref[prop:disall-factorized-subtraction]{10.5}, \bpBrk \hyperref[thm:disall-factorization]{11}；\bpBrk \hyperref[thm:appN-structure]{N.43} & 测试解除途径。 \\
\bpNumA{asm:appN-multren}\newline \bpLab{\detokenize{asm:}\bpBrk \detokenize{appN-}\bpBrk \detokenize{multren}} & \mbox{附录~N} & 共线因子的乘法重整化（维数正规化路线） (Z1)(Z3) & \bpT{RoyalBlue}{A} & 测试状态句。 & \hyperref[prop:pdf-timelike-dglap]{8.29}；\bpBrk \hyperref[thm:method-template]{9.42}；\bpBrk \hyperref[thm:disall-renorm-compat]{10.8}, \bpBrk \hyperref[prop:disall-h5]{9}, \bpBrk \hyperref[thm:disall-factorization]{11}；\bpBrk \hyperref[thm:coll-sia]{11.5}；\bpBrk \hyperref[asm:appN-ir-regulator]{\textcolor{orange!65!black}{N.33}}, \bpBrk \hyperref[thm:appN-structure]{43} & 测试解除途径。 \\
\bpNumA{asm:appN-ir-regulator}\newline \bpLab{\detokenize{asm:}\bpBrk \detokenize{appN-}\bpBrk \detokenize{ir-}\bpBrk \detokenize{regulator}} & \mbox{附录~N} & 保持 Ward 恒等式的部分子靶红外正规化 (Z2) & \bpT{Plum}{C} & 测试状态句。 & \hyperref[thm:disall-renorm-compat]{10.8} & 测试解除途径。 \\
\bpNumA{asm:appN-tmd}\newline \bpLab{\detokenize{asm:}\bpBrk \detokenize{appN-}\bpBrk \detokenize{tmd}} & \mbox{附录~N} & TMD 输入 (T1)--(T4) & \bpT{RoyalBlue}{A} & 测试状态句。 & \hyperref[prop:tmd-zero-bin-soft]{14.11}, \bpBrk \hyperref[thm:tmd-factorization]{15}, \bpBrk \hyperref[prop:tmd-evolution]{16}；\bpBrk \hyperref[prop:appN-tmd-window]{N.41}, \bpBrk \hyperref[thm:appN-structure]{43} & 测试解除途径。 \\
\bpNumA{asm:appN-endpoint}\newline \bpLab{\detokenize{asm:}\bpBrk \detokenize{appN-}\bpBrk \detokenize{endpoint}} & \mbox{附录~N} & 端点条件 (E1)--(E3) & \bpT{RoyalBlue}{A} & 测试状态句。 & \hyperref[prop:excl-soft-cancel]{19.3}, \bpBrk \hyperref[thm:excl-pion-ff]{4}, \bpBrk \hyperref[thm:excl-meson-L]{13}；\bpBrk \hyperref[thm:appN-structure]{N.43} & 测试解除途径。 \\
\bpNumA{asm:appN-analytic-reg}\newline \bpLab{\detokenize{asm:}\bpBrk \detokenize{appN-}\bpBrk \detokenize{analytic-}\bpBrk \detokenize{reg}} & \mbox{附录~N} & 两圈图样积分的辅助解析正规化 (RA1)(RA2) & \bpT{RoyalBlue}{A} & 测试状态句。 & \hyperref[prop:appN-twoloop-remainder]{N.51} & 测试解除途径。 \\
\bpNumA{asm:appO-lcopt}\newline \bpLab{\detokenize{asm:}\bpBrk \detokenize{appO-}\bpBrk \detokenize{lcopt}} & \mbox{附录~O} & 光锥序微扰论的适用性 & \bpT{RoyalBlue}{A} & 测试状态句。 & \hyperref[prop:method-fsi]{9.37}；\bpBrk \hyperref[lem:dy-fsi]{12.7}, \bpBrk \hyperref[thm:dy-factorization]{17}；\bpBrk \hyperref[prop:break-measurement]{18.2}；\bpBrk \hyperref[lem:appO-hard-locality]{O.11}, \bpBrk \hyperref[lem:appO-unitarity]{14}, \bpBrk \hyperref[prop:appO-ss-localize]{20}, \bpBrk \hyperref[lem:appO-side-independent-windows]{34}, \bpBrk \hyperref[prop:appO-ss-cancel]{35}, \bpBrk \hyperref[prop:appO-twoexch-spectator]{41}, \bpBrk \hyperref[prop:appO-one-sided]{54}, \bpBrk \hyperref[thm:appO-main]{68}--\hyperref[prop:appO-tmd]{70} & 测试解除途径。 \\
\bpNumA{asm:appO-adjacent-swap}\newline \bpLab{\detokenize{asm:}\bpBrk \detokenize{appO-}\bpBrk \detokenize{adjacent-}\bpBrk \detokenize{swap}} & \mbox{附录~O} & 相邻交换对称性（ADJ） (ADJ) & \bpT{BurntOrange}{P} & 测试状态句。 & \hyperref[prop:appO-twoexch-spectator]{O.41}, \bpBrk \hyperref[prop:appO-pergraph-weights]{44}, \bpBrk \hyperref[prop:appO-halfcoinc-sensitivity]{46} & 测试解除途径。 \\
\bpNumA{asm:appO-real-jets}\newline \bpLab{\detokenize{asm:}\bpBrk \detokenize{appO-}\bpBrk \detokenize{real-}\bpBrk \detokenize{jets}} & \mbox{附录~O} & 旁观者喷注函数的实性（REAL） (REAL) & \bpT{BurntOrange}{P} & 测试状态句。 & \hyperref[prop:appO-twoexch-spectator]{O.41}, \bpBrk \hyperref[prop:appO-pergraph-weights]{44}, \bpBrk \hyperref[prop:appO-halfcoinc-sensitivity]{46} & 测试解除途径。 \\
\bpNumA{asm:appO-collinear-ends}\newline \bpLab{\detokenize{asm:}\bpBrk \detokenize{appO-}\bpBrk \detokenize{collinear-}\bpBrk \detokenize{ends}} & \mbox{附录~O} & 共线端点的匹配 & \bpT{RoyalBlue}{A} & 测试状态句。 & \hyperref[prop:appO-twoexch-spectator]{O.41}, \bpBrk \hyperref[prop:appO-twoexch-regulated]{42}, \bpBrk \hyperref[prop:appO-halfcoinc-sensitivity]{46} & 测试解除途径。 \\
\bpNumA{asm:appO-envelope}\newline \bpLab{\detokenize{asm:}\bpBrk \detokenize{appO-}\bpBrk \detokenize{envelope}} & \mbox{附录~O} & 纤维积分误差的可积包络 & \bpT{RoyalBlue}{A} & 测试状态句。 & \hyperref[prop:method-glauber-deformation]{9.40}；\bpBrk \hyperref[prop:dy-deformation]{12.9}, \bpBrk \hyperref[thm:dy-factorization]{17}；\bpBrk \hyperref[asm:appO-mixing]{\textcolor{orange!65!black}{O.53}}, \bpBrk \hyperref[prop:appO-one-sided]{54}, \bpBrk \hyperref[thm:appO-main]{68}--\hyperref[prop:appO-tmd]{70} & 测试解除途径。 \\
\bpNumA{asm:appO-mixing}\newline \bpLab{\detokenize{asm:}\bpBrk \detokenize{appO-}\bpBrk \detokenize{mixing}} & \mbox{附录~O} & Glauber–软混合与重合连接的完整图集相消 & \bpT{RoyalBlue}{A} & 测试状态句。 & \hyperref[prop:method-glauber-deformation]{9.40}；\bpBrk \hyperref[prop:dy-deformation]{12.9}, \bpBrk \hyperref[thm:dy-factorization]{17}；\bpBrk \hyperref[lem:tmd-glauber]{14.13}, \bpBrk \hyperref[thm:tmd-factorization]{15}；\bpBrk \hyperref[prop:appO-one-sided]{O.54}, \bpBrk \hyperref[thm:appO-main]{68}--\hyperref[prop:appO-tmd]{70} & 测试解除途径。 \\
\bpNumA{asm:appO-deformation}\newline \bpLab{\detokenize{asm:}\bpBrk \detokenize{appO-}\bpBrk \detokenize{deformation}} & \mbox{附录~O} & 一般构型、相容性、软角区与 H-发射端点 (CS1)(CS2) & \bpT{RoyalBlue}{A} & 测试状态句。 & \hyperref[prop:method-glauber-deformation]{9.40}；\bpBrk \hyperref[prop:dy-deformation]{12.9}, \bpBrk \hyperref[thm:dy-factorization]{17}；\bpBrk \hyperref[lem:tmd-glauber]{14.13}；\bpBrk \hyperref[prop:appN-dy]{N.37}；\bpBrk \hyperref[thm:appO-main]{O.68}--\hyperref[prop:appO-tmd]{70} & 测试解除途径。 \\
\bpNumA{asm:appO-rewritten-cs}\newline \bpLab{\detokenize{asm:}\bpBrk \detokenize{appO-}\bpBrk \detokenize{rewritten-}\bpBrk \detokenize{cs}} & \mbox{附录~O} & 改写形式与族求和形式的 Collins 条件 (CS1)(CS2) & \bpT{RoyalBlue}{A} & 测试状态句。 & \hyperref[prop:method-glauber-deformation]{9.40}；\bpBrk \hyperref[prop:dy-deformation]{12.9}, \bpBrk \hyperref[thm:dy-factorization]{17}；\bpBrk \hyperref[lem:tmd-glauber]{14.13}, \bpBrk \hyperref[thm:tmd-factorization]{15}；\bpBrk \hyperref[prop:appN-dy]{N.37}；\bpBrk \hyperref[prop:appO-past-wilson]{O.65}, \bpBrk \hyperref[thm:appO-main]{68}--\hyperref[prop:appO-tmd]{70} & 测试解除途径。 \\
\end{longtable}
\endgroup
